\chapter{Objetivos}
\label{ch:objetivos}

En este capítulo se establecen las metas que guían el trabajo y sirven de referencia para evaluar sus resultados. Primero se enuncia el objetivo general, que fija el propósito global del proyecto; después se desglosa en una serie de objetivos específicos formulados según el criterio \gls{SMART} ---simples, medibles, acordados con la dirección, realistas y temporizados--- y ordenados de forma incremental, desde la adquisición de los fundamentos hasta la integración final del sistema. Esta jerarquía de objetivos actúa como hilo conductor del resto de la memoria.

\section{Objetivo general}

El objetivo general de este trabajo es el diseño e implementación de un sistema de planificación y control de maniobras orbitales que combine \acrfull{RL} con la capacidad de razonamiento en lenguaje natural de un \gls{LLM}, sobre un entorno de simulación física validado. El sistema se concibe con una arquitectura híbrida: agentes de \gls{RL} que aprenden las maniobras donde no existe una solución analítica sencilla, solucionadores clásicos donde el óptimo ya se conoce, y una capa \gls{LLM}  que orquesta esas herramientas y explica sus resultados. Con ello se persigue, por un lado, optimizar la toma de decisiones en escenarios de simulación espacial y, por otro, acercar la planificación de trayectorias a usuarios no especializados mediante una interacción en lenguaje natural.

\section{Objetivos específicos}

Para alcanzar el objetivo general, se definen los siguientes objetivos específicos, formulados según el criterio \gls{SMART} (simples, medibles, acordados con la dirección del trabajo, realistas y temporizados). Se enuncian de menor a mayor dependencia: cada uno se apoya en el anterior, de modo que el sistema se construye de forma incremental.

\begin{enumerate}

  \item \textbf{Adquirir y consolidar las competencias en mecánica orbital y en el manejo de la biblioteca Poliastro}. Antes de abordar el sistema, se afianzan los fundamentos de mecánica orbital y el dominio de la herramienta mediante la resolución de ejercicios prácticos de complejidad creciente: cálculo y visualización de órbitas, impulsos y maniobras, transferencias de Hohmann y trayectorias interplanetarias. Esta fase establece la base técnica y verifica que las herramientas se manejan con soltura antes de construir sobre ellas. La medición se realiza en base a completar un conjunto de ejercicios que cubra los tipos de maniobra y visualización empleados después en el proyecto, con resultados contrastables frente a valores conocidos (por ejemplo, el $\Delta v$ de una transferencia de Hohmann de referencia). Los límites son que esta es una fase de capacitación y no de producción de resultados originales; se apoya en la documentación de Poliastro y en la bibliografía de astrodinámica y la temporización de este objetivo se centra en la fase inicial del proyecto, previa al desarrollo del entorno. Se parte desde los conocimientos de mecánica orbital adquiridos en el grado.

  \item \textbf{Construir y validar un entorno de simulación física de mecánica orbital.} Se desarrolla, sobre la biblioteca Poliastro, un entorno capaz de modelar órbitas, maniobras impulsivas y la propagación con perturbaciones (achatamiento $J_2$ y arrastre atmosférico), incluyendo un modelo atmosférico propio de los cuerpos del Sistema Solar. Este entorno es el marco físico sobre el que el resto del sistema actúa, por lo que su fidelidad es requisito previo de todo lo demás. Para la medición, el modelo atmosférico reproduce los datos de las misiones de referencia con un factor de desviación inferior a 2, y la propagación coincide con la teoría analítica clásica (regresión nodal de Brouwer\footnote{Dirk Brouwer formuló una teoría analítica clásica que predice cómo el achatamiento del planeta ($J_2$) hace girar lentamente el plano de la órbita; se emplea aquí como referencia para validar el propagador.}) con un error inferior al 1\,\% en cuerpos de $J_2$ moderado; se acompaña de una batería de pruebas automáticas que debe superarse en su totalidad. Los límites se modelan nueve cuerpos (siete con atmósfera); no se persigue un modelo atmosférico de alta resolución temporal ni dependiente de la latitud, sino un perfil medio suficiente para el arrastre orbital y la temporización se centra en la fase inicial del proyecto. Se parte de las competencias y los \textit{scripts}
  desarrollados en el objetivo~1.

  \item \textbf{Desarrollar agentes de \gls{RL} para transferencias impulsivas, validados frente al óptimo analítico.} Se diseña un agente que aprende, sin conocer la fórmula, la transferencia entre órbitas, comparándolo con el óptimo de Hohmann que actúa de juez. El objetivo se descompone en tres subobjetivos de dificultad creciente:
  
  \begin{enumerate}
    \item un agente especializado en un caso fijo (de \gls{LEO} a \gls{GEO}), como prueba de concepto.
    \item un agente general que, mediante la adimensionalización del problema, resuelve cualquier transferencia coplanaria en cualquier cuerpo sin reentrenar;
    \item una extensión tridimensional que incorpora el cambio de plano.
  \end{enumerate}
  
  La medición se centra en el coste $\Delta v$ del agente que se desvía menos de un 5\,\% del óptimo de Hohmann (criterio de éxito acordado), alcanzando la órbita objetivo con un error inferior al 2\,\%. Los límites son las maniobras impulsivas entre órbitas circulares; la extensión 3D se restringe a las transferencias de subida con cambio de plano, que es el caso de interés práctico y la temporización se centra en la fase intermedia del proyecto.

  \item \textbf{Desarrollar un agente de \gls{RL} para una maniobra sin solución analítica cerrada: el aerofrenado multi-planeta.} Se diseña un agente que rebaja una órbita aprovechando el arrastre atmosférico en sucesivas pasadas por el perigeo ---usando el modelo atmosférico validado del objetivo~2---, controlando el riesgo mediante un criterio de presión dinámica. Es la maniobra donde el \gls{RL} resulta verdaderamente imprescindible, al no existir una fórmula que la resuelva, para ello la medición se realiza en base a que el agente alcance la órbita objetivo sin superar el límite de seguridad en una batería de escenarios aleatorios, y mejora de forma cuantificable (en número de pasadas) frente a una estrategia ingenua de referencia. Los límites de esto es que se entrena un especialista por planeta, dado que cada atmósfera tiene su propia escala; se excluyen los cuerpos sin atmósfera y aquellos en los que el aerofrenado es físicamente inviable y la temporización consta de la fase intermedia--final del proyecto.

\item \textbf{Desarrollar un agente de RL para el mantenimiento orbital
  (\textit{station-keeping}).} Se diseña un agente que conserva una órbita operativa a lo largo del tiempo, compensando con impulsos de re-boost la degradación por arrastre atmosférico. A diferencia de las maniobras anteriores ---de uno o dos impulsos---, es un problema temporal de control en lazo cerrado, en el que el agente decide cuándo y cuánto corregir a lo largo de toda la misión. La medición se realiza en base a que el agente mantenga la órbita dentro de una banda de tolerancia durante toda la misión en una batería de escenarios aleatorios y en todos los cuerpos con atmósfera, con un coste de combustible comparable o inferior al de una estrategia de referencia. Se corrige la geometría de la órbita (principalmente el semieje) frente al arrastre; la precesión del plano por $J_2$ se modela pero no se combate, por tratarse de una corrección fuera de plano de coste prohibitivo; se excluyen los cuerpos sin atmósfera, donde la órbita no se degrada pudiendo incluir esos casos como límites del agente y la temporización consta de la fase final del proyecto, además el punto de partida es el modelo físico del objetivo~2 y la experiencia de RL de los objetivos~3 y~4.

  \item \textbf{Integrar una capa de orquestación y explicación basada en un \gls{LLM}.} Se implementa un demostrador en el que un \gls{LLM} interpreta peticiones en lenguaje natural, selecciona y ejecuta la herramienta adecuada (un agente de \gls{RL} o un solucionador clásico) y explica el resultado, todo ello anclado a las cifras reales que devuelven las herramientas. El \gls{LLM}  no realiza cálculos orbitales ni se reentrena: su función es coordinar y comunicar (patrón de \textit{tool use}). La medición consta de que el sistema resuelva correctamente un conjunto representativo de peticiones en lenguaje natural, invocando en cada caso la herramienta correcta y presentando únicamente datos procedentes de la simulación (sin información inventada), los límites que surgen aquí es que se trata de un demostrador de la arquitectura, no de un producto; su alcance está condicionado a la disponibilidad de recursos de cómputo (un servicio de \gls{LLM}  por \gls{API}  o un modelo local en un servidor con capacidad suficiente), un condicionante que finalmente se resolvió mediante un servicio de LLM en la nube. La temporización consta de la fase final del proyecto.

  \item \textbf{Evaluar el sistema mediante métricas de ingeniería aeroespacial.} Se realiza una evaluación comparativa que cuantifica el rendimiento de los agentes frente a las referencias clásicas, empleando métricas como el coste $\Delta v$, el error de llegada, el tiempo de transferencia o el número de pasadas de aerofrenado, y se documentan con honestidad las limitaciones y los supuestos del modelo. La medición consta de disponer, para cada maniobra, de la comparación numérica agente--referencia y de las limitaciones identificadas y la temporización transversal, a lo largo de todo el periodo de trabajo.

\end{enumerate}
